% Slides — Capítulo 7: Processos de Precipitação
\documentclass[aspectratio=169,11pt]{beamer}
\usepackage{../comum/temameteo}
\graphicspath{{../figuras/cap07/}}

\title[Cap. 7 --- Precipitação]{Capítulo 7 --- Processos de Precipitação}
\subtitle{Meteorologia Prática}
\author{}
\date{}

\begin{document}

\begin{frame}
  \titlepage
  \vfill
  \atribuicao
\end{frame}

\begin{frame}{Roteiro}
  \tableofcontents
\end{frame}

% ---------------------------------------------------------------
\section{O problema das seis ordens de grandeza}

\begin{frame}{Da poeira à gota de chuva}
  \eqdestaque{$\underbrace{0{,}1\,\mu\mathrm{m}}_{\text{CCN}} \to
    \underbrace{10\,\mu\mathrm{m}}_{\text{gotícula}} \to
    \underbrace{1\,\mathrm{mm}}_{\text{gota de chuva}}$\qquad
    volume $\times10^6$ em $\sim$30\,min}
  \vspace{0.5em}
  \begin{itemize}
    \item Condensou (Cap.~6) ainda não é choveu.
    \item Difusão de vapor constrói a \alert{nuvem}; a \alert{chuva}
          precisa de colisões (quente) ou gelo (Bergeron).
    \item No fim: a distribuição de gotas que o radar enxerga (Cap.~8).
  \end{itemize}
\end{frame}

% ---------------------------------------------------------------
\section{Nucleação: Kelvin e Köhler}

\begin{frame}{Sem aerossol não há nuvem}
  \begin{columns}
    \column{0.4\textwidth}
    \eqdestaque{$\dfrac{e_s(r)}{e_s(\infty)} = e^{A/r}$ (Kelvin)}
    \vspace{0.4em}
    \begin{itemize}
      \item Gota de $0{,}01\,\mu$m exige UR $\simeq112\%$; a atmosfera
            para em $\sim100{,}5\%$.
      \item Nucleação homogênea: proibida. Toda nuvem nasceu sobre um
            grão (CCN).
    \end{itemize}
    \column{0.3\textwidth}
    \figlivroslide{fig_07_4b.jpg}{0.4\textheight}{Fig.~7.4 --- nucleação
      heterogênea sobre aerossol molhável}
    \column{0.3\textwidth}
    \figlivroslide{fig_07_6.jpg}{0.4\textheight}{Fig.~7.6 --- barreira de
      Kelvin: evaporar × crescer}
  \end{columns}
\end{frame}

\begin{frame}{Curvas de Köhler: Kelvin × Raoult}
  \begin{columns}
    \column{0.4\textwidth}
    \eqdestaque{$S(r) = 1 + \dfrac{A}{r} - \dfrac{B r_d^3}{r^3}$}
    \vspace{0.4em}
    \begin{itemize}
      \item Sal reduz $e_s$ (Raoult) e enfrenta a curvatura (Kelvin).
      \item Máximo $=$ ($r^*$, $S^*$): passou, \alert{ativou} — cresce
            sozinha.
      \item Núcleo maior/mais salgado ativa antes.
    \end{itemize}
    \column{0.6\textwidth}
    \figlivroslide{fig_07_7a.jpg}{0.52\textheight}{Fig.~7.7 --- curvas de
      Köhler para vários solutos}
  \end{columns}
\end{frame}

\begin{frame}{Ativação e o elo aerossol--clima}
  \begin{columns}
    \column{0.4\textwidth}
    \begin{itemize}
      \item Subida mais rápida $\Rightarrow$ $S_{max}$ maior
            $\Rightarrow$ mais núcleos ativados (\texttt{pyrcel},
            Caso~2).
      \item Marítimo limpo: $\sim10^2$ gotas/cm$^3$; continental
            poluído: $\sim10^3$ — nuvens mais brancas (Cap.~21).
      \item Gelo: nucleação exige IN raros ($1/10^5$ dos CCN); água
            super-resfriada até $-38\,°$C.
    \end{itemize}
    \column{0.6\textwidth}
    \figlivroslide{fig_07_8a.jpg}{0.52\textheight}{Fig.~7.8 --- o caminho
      de ativação de uma gota pela sua curva de Köhler}
  \end{columns}
\end{frame}

% ---------------------------------------------------------------
\section{Crescimento: difusão, coleta, Bergeron}

\begin{frame}{Difusão: a raiz quadrada que não chega lá}
  \begin{columns}
    \column{0.44\textwidth}
    \eqdestaque{$r\dfrac{dr}{dt} = C \Rightarrow r \propto \sqrt{t}$}
    \vspace{0.4em}
    \begin{itemize}
      \item 1$\to$10\,$\mu$m: minutos. 10\,$\mu$m$\to$1\,mm:
            \alert{dias}.
      \item Cada década de raio custa $100\times$ mais tempo.
      \item A difusão faz a nuvem; \alert{não} faz a chuva (Caso~3).
    \end{itemize}
    \column{0.56\textwidth}
    \figlivroslide{fig_07_12a.jpg}{0.5\textheight}{Fig.~7.12 ---
      crescimento difusivo: $R(t) \propto \sqrt{t}$}
  \end{columns}
\end{frame}

\begin{frame}{Colisão--coalescência: a avalanche da chuva quente}
  \begin{columns}
    \column{0.4\textwidth}
    \eqdestaque{$\dfrac{dr}{dt} = \dfrac{E\,\bar L\,\Delta w_T}{4\rho_L}$}
    \vspace{0.4em}
    \begin{itemize}
      \item Gota grande cai mais rápido e varre as pequenas — cresce
            \alert{acelerando}.
      \item Cruzamento difusão$\to$coleta em $\sim20\,\mu$m: o gargalo
            (N3).
      \item Nuvem poluída (muitas gotas pequenas) demora a
            atravessá-lo.
    \end{itemize}
    \column{0.3\textwidth}
    \figlivroslide{fig_07_19.jpg}{0.42\textheight}{Fig.~7.19 ---
      eficiência de colisão: trajetórias em volta do coletor}
    \column{0.3\textwidth}
    \figlivroslide{fig_07_20.jpg}{0.42\textheight}{Fig.~7.20 ---
      coalescência eficiente × ineficiente}
  \end{columns}
\end{frame}

\begin{frame}{Bergeron: o gelo rouba das gotículas}
  \begin{columns}
    \column{0.4\textwidth}
    \begin{itemize}
      \item $e_s^{gelo} < e_s^{agua}$: diferença máxima em
            $\alert{-12\,°C}$ (Cap.~4).
      \item O raro cristal vê ambiente supersaturado: cresce às custas
            das gotículas.
      \item Depois: agregação (flocos), riming (graupel),
            fragmentação.
      \item A chuva de latitudes médias nasce \alert{neve}.
    \end{itemize}
    \column{0.28\textwidth}
    \figlivroslide{fig_07_16.jpg}{0.4\textheight}{Fig.~7.16 ---
      $e_s^{agua}-e_s^{gelo}$: máximo em $-12\,°$C}
    \column{0.32\textwidth}
    \figlivroslide{fig_07_17.jpg}{0.4\textheight}{Fig.~7.17 ---
      agregação, riming e coalescência}
  \end{columns}
\end{frame}

\begin{frame}{Hábitos dos cristais}
  \begin{columns}
    \column{0.4\textwidth}
    \begin{itemize}
      \item Agulhas, placas, dendritos, colunas: o hábito depende de
            $T$ e do excesso de vapor.
      \item Dendritos ($\sim-15\,°$C): crescimento mais rápido — a zona
            de Bergeron por excelência.
      \item Cada floco conta a história térmica da sua queda.
    \end{itemize}
    \column{0.6\textwidth}
    \figlivroslide{fig_07_12c.jpg}{0.54\textheight}{Fig.~7.12 --- hábitos
      de cristal em função de temperatura e vapor}
  \end{columns}
\end{frame}

\begin{frame}{A fábrica completa dentro da nuvem fria}
  \begin{columns}
    \column{0.4\textwidth}
    \begin{itemize}
      \item Acima de $-40\,°$C: só gelo; $-40$ a $0\,°$C: mistura
            (Bergeron ativo); abaixo: derretimento.
      \item A fase que chega ao chão depende do perfil de $T_w$
            (Cap.~4): chuva, neve, chuva congelante.
      \item Versão convectiva e violenta: granizo (Cap.~14).
    \end{itemize}
    \column{0.6\textwidth}
    \figlivroslide{fig_07_21.jpg}{0.56\textheight}{Fig.~7.21 --- os
      andares de fase de uma nuvem precipitante fria}
  \end{columns}
\end{frame}

% ---------------------------------------------------------------
\section{Tamanhos, quedas e a ponte para o radar}

\begin{frame}{Velocidade terminal}
  \begin{columns}
    \column{0.44\textwidth}
    \eqdestaque{Stokes ($r<40\,\mu$m): $w_T = k_1 r^2$;\quad gotas de
      chuva: 4--9\,m/s}
    \vspace{0.4em}
    \begin{itemize}
      \item T4: peso $=$ arrasto viscoso $\Rightarrow$
            $w_T \propto r^2$.
      \item Gotas $>$ 3\,mm quebram: teto natural de tamanho.
      \item $\Delta w_T$ é o motor da coleta.
    \end{itemize}
    \column{0.56\textwidth}
    \figlivroslide{fig_07_18.jpg}{0.52\textheight}{Fig.~7.18 ---
      velocidade terminal vs.\ raio: os três regimes}
  \end{columns}
\end{frame}

\begin{frame}{Marshall--Palmer e a relação $Z$--$R$}
  \begin{columns}
    \column{0.44\textwidth}
    \eqdestaque{$N(D) = N_0\,e^{-\Lambda D}$,\quad
      $\Lambda = 4{,}1\,R^{-0{,}21}$\,mm$^{-1}$}
    \vspace{0.4em}
    \begin{itemize}
      \item Chuva forte $\Rightarrow$ $\Lambda$ menor $\Rightarrow$
            gotas grandes.
      \item $R \propto \int N w_T D^3$;\quad
            $Z \propto \int N D^6$ — o $D^6$ adora gotas grandes.
      \item N4 verifica $Z = aR^{1{,}4}$: a fórmula do radar (Cap.~8).
    \end{itemize}
    \column{0.56\textwidth}
    \figlivroslide{fig_07_23.jpg}{0.52\textheight}{Fig.~7.23 ---
      distribuição de tamanhos para chuvas de várias intensidades}
  \end{columns}
\end{frame}

\begin{frame}{Onde chove no planeta}
  \begin{columns}
    \column{0.4\textwidth}
    \begin{itemize}
      \item ITCZ: a faixa equatorial que veremos nascer no Cap.~11.
      \item Desertos subtropicais: o ramo que desce de Hadley.
      \item Trilhas de tempestades das latitudes médias (Cap.~13).
    \end{itemize}
    \column{0.6\textwidth}
    \figlivroslide{fig_07_25.jpg}{0.5\textheight}{Fig.~7.25 ---
      precipitação média anual global}
  \end{columns}
\end{frame}

% ---------------------------------------------------------------
\section{Síntese e laboratório}

\begin{frame}{O mapa do capítulo}
  \eqdestaque{CCN $\xrightarrow{\text{Köhler}}$ gotícula
    $\xrightarrow{\text{difusão }\sqrt{t}}$ $10\,\mu$m
    $\xrightarrow[\text{ou Bergeron}]{\text{colisões}}$ gota
    $\xrightarrow{\text{M--P}}$ radar}
  \vspace{0.6em}
  Sem aerossol não há nuvem; sem colisões ou gelo não há chuva. A
  distribuição final é exatamente o que o radar do Cap.~8 mede.
\end{frame}

\begin{frame}{Laboratório numérico --- notebook do Cap.~7}
  \texttt{notebooks/cap07\_precipitacao.ipynb}
  \begin{enumerate}
    \item Curvas de Köhler e ativação.
    \item Parcela subindo com \texttt{pyrcel}: $S_{max}$ e fração
          ativada.
    \item Difusão × coleta: o gargalo dos 20\,$\mu$m.
    \item Marshall--Palmer $\to$ $R$, $Z$ e a relação $Z$--$R$.
  \end{enumerate}
  \vspace{0.4em}
  \textbf{Exercícios}: T1--T5 e N1--N4 nas notas; células reservadas no
  notebook.
  \vfill
  \atribuicao
\end{frame}

\end{document}
